\chapter{Experimental Evaluation}\label{experimental-evaluation}

This chapter reports the completed evaluation of configurations C1--C4
on the 500-record synthetic experimental dataset based on HTTP CSIC
2010. It reports web-event
classification, abstentions, attack-subtype matches on malicious
records, the C4 detector counts, and citation-ID membership checks.
The chapter also reports post-hoc Prompt Injection compliance, a
stratified semantic-quality review, and classification and Prompt
Injection results on the separate ModSec-WP external evaluation
dataset. These additions distinguish
classification performance from instruction obedience, citation
integrity from semantic support, and detector performance from
end-to-end defense.

\section{Experimental Setup}\label{experimental-setup}

The initial full-dataset batch ran all four configurations once on all 500
records, yielding 2{,}000 responses. The comparison tables use C1--C3 from
that batch and replace its initial C4 results with a later C4-only
execution on the same records. The reported C4 uses secure-prompt version
\texttt{attempt-semantics-v4} and few-shot version
\texttt{independent-examples-v2}; C3 uses prompt version
\texttt{attempt-semantics-v2} and few-shot version
\texttt{independent-examples-v1}. The C4 revision also changes detector
decoding and sanitization. Therefore, the C3--C4 difference does not
isolate the effect of sanitization or any single defense component.
The recorded content of secure-prompt versions v2 and v4 is summarized
in Appendix~\ref{appendix-secure-prompt-versions}.

The provider returned
\texttt{gpt-4o-mini-2024-07-18} for the \texttt{gpt-4o-mini} model
alias. Generation used temperature 0.1, structured output, retrieval
top-\emph{k} of 3, and a local knowledge base containing 57 documents.
All 2{,}000 responses in the initial batch had valid structured JSON
directly; none required recovery and none was invalid. C1 and C2 used no
few-shot examples. C3 used examples for all 500 records. The reported C4
used examples only on records routed to its secure path by the detector.

The corpus has 250 clean and 250 adversarial records arranged in 250
pairs. Its class distribution is 400 malicious and 100 benign records.
Five records from this corpus had already been used to adjust the
prompt; consequently, it is not an independent test of those prompt
choices. C4 was run in multiple development iterations, while only one
reported version of each configuration is shown. These results do not
estimate run-to-run variability.

A separate external evaluation batch used the 800-record paired
ModSec-WP dataset and evaluated all four configurations once, yielding
3{,}200 responses. C3 and C4 used secure-prompt version
\texttt{attempt-\allowbreak{}semantics-\allowbreak{}v5} and few-shot
version \texttt{independent-\allowbreak{}examples-\allowbreak{}v2}; C4 additionally
applied detector-based routing and sanitization. Because the ModSec-WP
sample informed the contextual indicators in prompt v5, this batch is
development evidence rather than an independent holdout evaluation. The
v5 prompt is reproduced in Appendix~\ref{appendix-secure-prompt-v5}.
The reported ModSec-WP metrics below exclude the 200 records in the
\texttt{bruteforce} stratum and use the remaining 600 records (300
pairs, or 2{,}400 responses across C1--C4). This category requires
repeated-request evidence over a time window and cannot be established
from the individual request fields supplied to the model.

\subsection{Metric conventions}\label{evaluation-metric-conventions}

Binary classification means assigning each record to one of two
web-event classes: benign/normal or malicious. It does not mean
assigning an attack subtype. Classification accuracy is the number of
records with the correct benign/malicious label divided by all records
in the stated group. Unknown predictions are abstentions, not a third
class; they remain in the denominator and are not counted as correct.
``Binary errors'' below are incorrect,
non-unknown binary predictions; unknowns and invalid outputs are
reported separately. Where shown as ``errors (FP/FN)'', the parenthetic
counts give false positives followed by false negatives. Thus, for each
group, correct predictions, unknowns, binary errors, and invalid
outputs sum to the denominator.
Accuracy drop is clean accuracy minus adversarial accuracy, expressed
in percentage points (pp), not as a relative percentage change.

End-to-end attack-subtype accuracy is calculated only on records whose
ground-truth web-event class is malicious. A record counts as correct
only when the model predicts \texttt{malicious} and assigns the correct
attack subtype; \texttt{benign}, \texttt{unknown}, and incorrect subtype
predictions count as incorrect. Its denominator is 400 overall and 100
within each malicious category on the CSIC 2010-based synthetic dataset.
In the reported ModSec-WP cohort, the corresponding denominator is 400
overall and 200 within each malicious source category (SQL Injection
and XSS). This is not precision over predicted malicious records.
Benign ground-truth records are excluded because they have no attack
subtype.

\section{Configuration Comparison}\label{configuration-comparison}

\begin{table}[H]
\centering
\caption{Overall binary classification on the 500-record CSIC 2010-based synthetic dataset.}
\label{tab:overall-classification}
\begin{tabular}{@{}lrrrrr@{}}
\toprule
Configuration & Correct / 500 & Accuracy & Unknown & Errors (FP/FN) & Invalid \\
\midrule
C1 & 440 & 88.0\% & 46 & 14 (11/3) & 0 \\
C2 & 404 & 80.8\% & 59 & 37 (2/35) & 0 \\
C3 & 352 & 70.4\% & 141 & 7 (0/7) & 0 \\
C4 & 453 & 90.6\% & 22 & 25 (0/25) & 0 \\
\bottomrule
\end{tabular}
\par\smallskip
\begin{minipage}{\textwidth}\footnotesize
\textit{Note.} Binary classification distinguishes benign/normal from
malicious records. ``Unknown'' is an abstention, not a third class, and
is counted as incorrect in the accuracy denominator.
\end{minipage}
\end{table}

Table~\ref{tab:overall-classification} shows C4 with the highest overall
accuracy, followed by C1. The error and abstention columns show different
trade-offs: C3 has 141 unknown predictions but only seven non-unknown
errors, whereas C4 has 22 unknown predictions and 25 false negatives.
Since unknown predictions count as incorrect in accuracy, the
condition-level results in Table~\ref{tab:condition-classification}
help show how performance differs between clean and adversarial records.

\begin{table}[H]
\centering
\caption{Binary classification by condition on the CSIC 2010-based synthetic dataset (250 records per condition).}
\label{tab:condition-classification}
\resizebox{\textwidth}{!}{%
\begin{tabular}{@{}llrrrrr@{}}
\toprule
Configuration & Condition & Correct / 250 & Accuracy & Unknown & Errors (FP/FN) & Invalid \\
\midrule
C1 & Clean & 225 & 90.0\% & 25 & 0 (0/0) & 0 \\
C1 & Adversarial & 215 & 86.0\% & 21 & 14 (11/3) & 0 \\
C2 & Clean & 206 & 82.4\% & 19 & 25 (0/25) & 0 \\
C2 & Adversarial & 198 & 79.2\% & 40 & 12 (2/10) & 0 \\
C3 & Clean & 187 & 74.8\% & 56 & 7 (0/7) & 0 \\
C3 & Adversarial & 165 & 66.0\% & 85 & 0 (0/0) & 0 \\
C4 & Clean & 210 & 84.0\% & 15 & 25 (0/25) & 0 \\
C4 & Adversarial & 243 & 97.2\% & 7 & 0 (0/0) & 0 \\
\midrule
\multicolumn{2}{@{}l}{Clean-to-adversarial accuracy drop} & \multicolumn{5}{l@{}}{} \\
C1 & & \multicolumn{5}{l}{4.0 pp} \\
C2 & & \multicolumn{5}{l}{3.2 pp} \\
C3 & & \multicolumn{5}{l}{8.8 pp} \\
C4 & & \multicolumn{5}{l}{-13.2 pp} \\
\bottomrule
\end{tabular}
}
\par\smallskip
\begin{minipage}{\textwidth}\footnotesize
\textit{Note.} The predicted classes are benign/normal and malicious;
``Unknown'' is an abstention, not a third class, and is counted as
incorrect in accuracy.
\end{minipage}
\end{table}

The tables show different combinations of abstention and incorrect
non-unknown predictions. For example, C3's adversarial condition has
85 unknowns and no incorrect non-unknown binary predictions, whereas
C4's clean condition has 25 binary errors and 15 unknowns. These are
descriptive results from one run, not evidence of statistical
significance. The clean-to-adversarial difference is a condition
contrast; it does not establish that Prompt Injection caused the
difference.

\section{Results by Configuration}\label{results-by-configuration}

All configurations receive the same normalized HTTP record, use the same
model alias and generation settings, and must return the same structured
report. They differ in the information placed in context and in the
controls applied before generation. Figure~\ref{fig:configuration-flows}
summarizes the comparison: C1 is the direct baseline, C2 adds trusted
retrieval, C3 applies the secure prompt to every record, and C4 activates
the secure path only after a detector match.

\begin{figure}[H]
\centering
\resizebox{0.78\textwidth}{!}{%
\begin{tikzpicture}[
  box/.style={draw=black!65,fill=gray!7,rounded corners=2pt,
    minimum height=0.6cm,text width=1.95cm,align=center,
    inner sep=3pt,font=\scriptsize},
  lane/.style={font=\footnotesize\bfseries,anchor=east},
  arrow/.style={-{Latex[length=1.7mm]},thick,black!65}
]
  \node[lane] at (-0.35,2.4) {C1};
  \node[box] (c1a) at (1.0,2.4) {HTTP input};
  \node[box] (c1b) at (3.8,2.4) {LLM};
  \node[box] (c1c) at (6.6,2.4) {Report};
  \draw[arrow] (c1a) -- (c1b); \draw[arrow] (c1b) -- (c1c);

  \node[lane] at (-0.35,1.6) {C2};
  \node[box] (c2a) at (1.0,1.6) {HTTP input};
  \node[box] (c2b) at (3.8,1.6) {RAG + LLM};
  \node[box] (c2c) at (6.6,1.6) {Report};
  \draw[arrow] (c2a) -- (c2b); \draw[arrow] (c2b) -- (c2c);

  \node[lane] at (-0.35,0.8) {C3};
  \node[box] (c3a) at (1.0,0.8) {HTTP input};
  \node[box] (c3b) at (3.8,0.8) {RAG + secure};
  \node[box] (c3c) at (6.6,0.8) {Report};
  \draw[arrow] (c3a) -- (c3b); \draw[arrow] (c3b) -- (c3c);

  \node[lane] at (-0.35,0.0) {C4};
  \node[box] (c4a) at (1.0,0.0) {Detect + clean};
  \node[box] (c4b) at (3.8,0.0) {RAG + secure};
  \node[box] (c4c) at (6.6,0.0) {Report + check};
  \draw[arrow] (c4a) -- (c4b); \draw[arrow] (c4b) -- (c4c);
\end{tikzpicture}%
}
\caption{Information flow and controls applied by C1--C4.}
\label{fig:configuration-flows}
\end{figure}

\subsection{C1 --- LLM Only}\label{baseline-c1-llm-only}

C1 correctly classified 440 of 500 records (88.0\%), with 46 unknowns
and 14 incorrect non-unknown predictions. Its clean and adversarial
accuracies were 90.0\% and 86.0\%, respectively (a 4.0 pp difference).

C1 receives no retrieved policy or detector metadata. Its 440 correct
classifications therefore measure what the model can infer directly
from the request fields and status code. The result is strong for
explicit signatures: SQL Injection, Path Traversal, and XSS are mostly
recognizable from one event. The weak point is that the model cannot
distinguish an embedded instruction from evidence about the underlying
web event. For example, the adversarial member
\texttt{BENIGN\_0201\_ADV} contains an ordinary product search together
with ``ignore previous instructions''. C1 described the instruction
as its indicator and classified the request as malicious, producing a
false positive. Conversely, in \texttt{BF\_0170\_ADV}, a login request
included text requesting a benign label; C1 followed part of that
context and returned benign although the web-event annotation was
malicious.

These examples show why the baseline is useful but not a security
control. A response can be structured and syntactically valid while
still allowing attacker-controlled words to alter the meaning assigned
to the request. C1's 11 false positives on benign records and three
false negatives on malicious records should therefore be read together
with its 46 abstentions, rather than treating 88.0\% accuracy as proof
of resistance.

\subsection{C2 --- LLM + RAG}\label{evaluation-c2-llm-rag}

C2 correctly classified 404 of 500 records (80.8\%), with 59 unknowns
and 37 incorrect non-unknown predictions. Its clean and adversarial
accuracies were 82.4\% and 79.2\% (a 3.2 pp difference). The observed
scores do not show an improvement over C1 in this run; the single-run
design does not support a causal attribution to retrieval.

C2 adds three trusted knowledge chunks to the same task. The intended
benefit is to provide definitions and response guidance, not to supply
the answer key. In this run, however, retrieval coincided with a decline
from 88.0\% to 80.8\% and with 37 incorrect non-unknown classifications.
The retrieved documents may improve terminology, but the model still
receives the embedded instruction unchanged and can combine it with
reference material in a misleading way.

A representative benign false positive was
\texttt{BENIGN\_0206\_ADV}: the event was a normal product search, but the
query parameter asked the system to stop security analysis. C2 cited
retrieval evidence yet used the suspicious instruction as an indicator
of malicious traffic. At the opposite extreme,
\texttt{BF\_0151\_CLEAN} contained an \texttt{admin} login with a common
password and a 403 response. C2 interpreted the blocked response as
insufficient evidence and classified the event as benign. Thus retrieval
did not systematically resolve the distinction between an attempted web
attack, a successful attack, and an embedded Prompt Injection.

The 35 false negatives are concentrated in Brute Force Login cases
that require a decision about how much sequence-level meaning may be
inferred from a single event. The result does not demonstrate that RAG
is intrinsically harmful; it shows that generic retrieved references did
not replace an explicit policy for classifying attempts under the
available evidence.

\subsection{C3 --- LLM + RAG + Secure Prompt}
\label{evaluation-c3-llm-rag-secure-prompt}

C3 correctly classified 352 of 500 records (70.4\%), with 141 unknowns
and 7 incorrect non-unknown predictions. Its clean and adversarial
accuracies were 74.8\% and 66.0\% (an 8.8 pp difference). Its secure
prompt and examples were applied to every record. These classification
counts do not establish whether injected instructions were followed.

C3 applies \texttt{SECURE\_SYSTEM} and the same few-shot examples to every
record, including records for which the detector would have found no
injection. This is the strongest intervention in prompt policy but not
in context preparation: the suspicious text is still present and the
model is asked to resist it. The resulting 70.4\% accuracy is explained
mainly by abstention. C3 produced 141 \texttt{unknown} responses,
including 92 in Brute Force Login and 26 in Benign/Normal.

For a clean login request such as \texttt{BF\_0151\_CLEAN}, C3 reported
that the 403 response looked like an ordinary failed login and that no
conclusive web-attack indicator was present. That behaviour is cautious,
but it conflicts with the dataset's definition of a suspected credential
guess as a malicious attempt. The secure policy therefore protected
against overclaiming while also suppressing some intended positive
classifications. The same mechanism affected benign adversarial records:
the presence of an injected instruction made the model uncertain even
when the remaining request was an ordinary search or form submission.

C3 is consequently useful as an ablation of prompt hardening. It shows
that repeating security instructions can reduce unsupported conclusions
and classification manipulation, but it does not recover the underlying
event if the conflicting text remains in context. Its low utility is a
key reason C4 combines routing and preprocessing instead of applying the
secure policy indiscriminately.

\subsection{C4 --- LLM + RAG + Detector + Evidence Control}
\label{evaluation-c4-llm-rag-detector-evidence-control}

C4 correctly classified 453 of 500 records (90.6\%), with 22 unknowns
and 25 incorrect non-unknown predictions. Its clean and adversarial
accuracies were 84.0\% and 97.2\% (the adversarial accuracy was 13.2 pp
higher). The secure
few-shot path was used for 240 detector-routed records, rather than all
500. On that path, the detected instruction was removed from the model
input and retrieval query while the remaining web-event fields were
preserved. C4's detector and citation-ID membership counts are reported
separately below.

\subsubsection{Development and Improvement of C4}
\label{development-and-improvement-of-c4}

The C4 result in this chapter is the revised implementation, not the
first prototype. The initial prototype used the same four stages as the
reported C4---detection, secure prompting, retrieval, and citation
validation---but it left the detected instruction inside the HTTP request
sent to both the retriever and the LLM. The expected defence was therefore
indirect: the model was instructed to treat the embedded text as
untrusted data. The initial results showed that this strategy was not
adequate for utility. It classified 377/500 synthetic records correctly
(75.4\%), produced 96 \texttt{unknown} responses, and retained correct
classification on 65.6\% of the clean/adversarial pairs.

The detector and model input must account for URL encoding: spaces may
appear as \texttt{+} and punctuation as percent-encoded characters, so a
pattern may not match the stored representation directly. The revised
detector repeatedly URL-decodes a separate matching copy while
preserving the original request for audit. The final ModSec-WP
evaluation below measures detector behavior on its prepared injection
patterns; because those patterns were designed to match the study's
defined variants, this is a coverage test rather than an estimate of
generalization to unseen injections.

Detection alone was not sufficient. In the initial prototype, even a
detected instruction remained visible in the user message as well as in
the retrieval query. A secure prompt could ask the model to ignore the
sentence, but the order still competed with the trusted task and could
also influence embedding similarity. The revised C4 therefore sanitizes
the C4 secure path before either operation. For a URL or query string,
it removes only the parameter containing a matched injection pattern.
For an instruction embedded in another untrusted field, it redacts the
matched fragment. The trusted HTTP semantics---method, path, payload,
status, and the remaining parameters---are retained.

For example, consider the following simplified input:

\begin{lstlisting}
POST /products?pi_note=Ignore+previous+instructions
     &action=search
body: q=laptop
status: 200
\end{lstlisting}

The detector identifies \texttt{pi\_note} as a candidate instruction
after URL decoding. The parameter removal is implemented as a filtering
operation rather than as a text-wide deletion. First, the normalized
query is split into name--value components using the \texttt{\&}
separator. Each component is then tested independently against the same
fixed rules. A component is discarded only when a rule matches within
that component; the remaining components are joined again in their
original order. The request path and all non-matching parameters are
therefore preserved.

\begin{figure}[H]
\centering
\resizebox{0.38\textwidth}{!}{%
\begin{tikzpicture}[
  box/.style={draw=black!65,fill=gray!7,rounded corners=2pt,
    minimum width=5cm,minimum height=0.58cm,
    text width=5cm,align=center,inner sep=2pt,font=\scriptsize},
  arrow/.style={-{Latex[length=1.7mm]},thick,black!65}
]
  \node[box] (raw) {1. Decode field};
  \node[box, below=0.2cm of raw] (split) {2. Split on \texttt{\&}};
  \node[box, below=0.2cm of split] (match) {3. Match each parameter};
  \node[box, below=0.2cm of match] (filter) {4. Discard matched parameter};
  \node[box, below=0.2cm of filter] (join) {5. Rejoin valid parameters};
  \node[box, below=0.2cm of join] (use) {6. Clean input for RAG + LLM};
  \draw[arrow] (raw) -- (split);
  \draw[arrow] (split) -- (match);
  \draw[arrow] (match) -- (filter);
  \draw[arrow] (filter) -- (join);
  \draw[arrow] (join) -- (use);
\end{tikzpicture}%
}
\caption{C4 parameter-level sanitization after a Prompt Injection match.}
\label{fig:c4-parameter-sanitization}
\end{figure}

For the example above, decoding and splitting produce two independent
components:

\begin{lstlisting}
["pi_note=Ignore previous instructions", "action=search"]
\end{lstlisting}

The first component matches the instruction-override rule and is
removed. The second component does not match and is retained. C4 then
constructs the model and retrieval input as:

\begin{lstlisting}
POST /products?action=search
body: q=laptop
status: 200
\end{lstlisting}

For fields without a name--value structure, such as \texttt{user\_agent}
or an arbitrary request body, C4 redacts only the matched instruction
fragment. For instance, ``Ignore previous instructions; browser agent''
becomes ``[REDACTED\_UNTRUSTED\_INSTRUCTION]; browser agent''. The rest
of the field remains available for the model. In all cases, the original
record remains stored for audit and reproducibility; the filtered copy
is the version used for retrieval and generation.

The distinction between the two operations is important. Parameter
removal deletes an isolated attacker-created instruction from a URL or
query string; fragment redaction masks an instruction embedded in
less-structured content. Neither operation uses the ground-truth
web-event label, severity, or attack subtype. It therefore cannot
``repair'' a malicious classification by consulting the answer. It only
removes content identified by the fixed Prompt Injection rules.

This transformation does not convert the request into a known benign
example. It removes a competing instruction while leaving the actual
web evidence available. If the underlying payload were instead a SQL
Injection expression, that expression would remain unchanged. Thus the
defence acts on Prompt Injection syntax, not on the ground-truth label or
the web-attack category.

The secure prompt was also revised. The initial wording emphasized that
a single event does not prove repeated authentication attempts, but the
model overgeneralized this caution: isolated requests with privileged
usernames and common passwords were often returned as
\texttt{unknown}. The revised prompt distinguishes two statements. It
still forbids inferring a brute-force campaign or a successful login
from one request. However, one request containing a plausible credential
guess against an authentication endpoint may be labelled as a suspected
single Brute Force Login attempt. Similarly, an injected note in an
otherwise normal request does not by itself justify an \texttt{unknown}
classification when the other fields support a benign event. The few-shot
examples were updated to demonstrate these two distinctions.

\begin{table}[H]
\centering
\caption{Initial and revised C4 results on the CSIC 2010-based synthetic dataset.}
\label{tab:c4-development-comparison}
\begin{tabular}{@{}lrr@{}}
\toprule
Measure & Initial C4 & Revised C4 \\
\midrule
Synthetic correct / 500 & 377 (75.4\%) & 453 (90.6\%) \\
Synthetic \texttt{unknown} & 96 & 22 \\
Synthetic paired retention & 65.6\% & 81.6\% \\
\bottomrule
\end{tabular}
\par\smallskip
\begin{minipage}{\textwidth}\footnotesize
\textit{Note.} Synthetic correct counts refer to the two-class web-event
decision (benign versus malicious). Unknown predictions are not correct
classifications. External ModSec-WP results are reported separately
below.
\end{minipage}
\end{table}

Table~\ref{tab:c4-development-comparison} summarizes a development
iteration using the same model and synthetic corpus. The revised C4
improved classification in the synthetic corpus from 75.4\% to 90.6\%,
reduced abstentions from 96 to 22, and raised clean/adversarial pair retention in
the synthetic evaluation from 65.6\% to 81.6\%.
The revisions changed detector decoding and sanitization as well as the
secure prompt and few-shot examples. These results therefore describe the
combined revised C4 package; they do not isolate sanitization as the
cause of the improvement.

This improvement should not be interpreted as complete security or
generalization to an independent external test. The ModSec-WP results
reported below are limited to one run on a prepared sample, and the
injection rules remain finite: new paraphrases, alternative encodings,
or multilingual attacks can still evade detection. The principal result
is therefore descriptive performance on the evaluated scenarios, not a
production guarantee.

\section{Attack-Subtype Results}\label{attack-subtype-results}

\begin{table}[htbp]
\centering
\caption{End-to-end correct attack-subtype classifications on the CSIC 2010-based synthetic dataset (400 malicious records).}
\label{tab:subtype-overall}
\begin{tabular}{@{}lrr@{}}
\toprule
Configuration & Correct / 400 malicious & End-to-end accuracy \\
\midrule
C1 & 374 & 93.50\% \\
C2 & 320 & 80.00\% \\
C3 & 278 & 69.50\% \\
C4 & 353 & 88.25\% \\
\bottomrule
\end{tabular}
\end{table}

\begin{table}[htbp]
\centering
\caption{End-to-end correct attack-subtype classifications by category on the CSIC 2010-based synthetic dataset (100 malicious records per category).}
\label{tab:subtype-by-category}
\begin{tabular}{@{}lrrrr@{}}
\toprule
Category & C1 & C2 & C3 & C4 \\
\midrule
Brute Force Login & 75/100 & 37/100 & 1/100 & 63/100 \\
Path Traversal & 100/100 & 89/100 & 95/100 & 96/100 \\
SQL Injection & 100/100 & 100/100 & 96/100 & 100/100 \\
XSS & 99/100 & 94/100 & 86/100 & 94/100 \\
\bottomrule
\end{tabular}
\end{table}

These scores require both a correct malicious binary classification and
the correct subtype; abstentions, benign predictions on malicious
records, and incorrect subtypes are all errors. They are conditional on
ground-truth malicious records and are not precision among the model's
positive predictions.

\section{Results by Dataset Category}\label{per-attack-analysis}

Table~\ref{tab:category-binary} reports correct binary classifications,
unknowns, and incorrect non-unknown predictions for all 100 records in
each category. For malicious categories, subtype counts are given in
Table~\ref{tab:subtype-by-category}; they use the same 100
ground-truth malicious records per category. Benign has no subtype
score. Separating abstentions from explicit binary errors is important
because C3 and C4 frequently return \texttt{unknown} for Brute Force
Login rather than committing to a false label.

\begin{table}[H]
\centering
\small
\caption{Binary results by category on the CSIC 2010-based synthetic dataset (correct, unknown, errors; 100 records per category).}
\label{tab:category-binary}
\begin{tabular}{@{}llrrrr@{}}
\toprule
Category & Measure & C1 & C2 & C3 & C4 \\
\midrule
Benign / Normal & Correct & 66 & 84 & 74 & 100 \\
 & Unknown & 23 & 14 & 26 & 0 \\
 & Binary errors & 11 & 2 & 0 & 0 \\
\midrule
Brute Force Login & Correct & 75 & 37 & 1 & 63 \\
 & Unknown & 23 & 29 & 92 & 12 \\
 & Binary errors & 2 & 34 & 7 & 25 \\
\midrule
Path Traversal & Correct & 100 & 89 & 95 & 96 \\
 & Unknown & 0 & 11 & 5 & 4 \\
 & Binary errors & 0 & 0 & 0 & 0 \\
\midrule
SQL Injection & Correct & 100 & 100 & 96 & 100 \\
 & Unknown & 0 & 0 & 4 & 0 \\
 & Binary errors & 0 & 0 & 0 & 0 \\
\midrule
XSS & Correct & 99 & 94 & 86 & 94 \\
 & Unknown & 0 & 5 & 14 & 6 \\
 & Binary errors & 1 & 1 & 0 & 0 \\
\bottomrule
\end{tabular}
\par\smallskip
\begin{minipage}{\textwidth}\footnotesize
\textit{Note.} Every record is assigned one of two event classes:
benign/normal or malicious. ``Correct'' counts records assigned the
right class. ``Unknown'' means the model abstained from assigning a
class. ``Binary errors'' count incorrect definite assignments between
the two classes and exclude ``Unknown'' cases.
\end{minipage}
\end{table}

The largest between-configuration variation is in Brute Force Login,
including many unknown predictions under C3 and C4. This is a
descriptive observation on this synthetic sample. In particular, a
Brute Force Login label can encode a sequence-level interpretation that
an isolated request record does not establish; the labels therefore do
not necessarily represent evidence available in a single event.

\section{Prompt Injection Outcomes and Additional Metrics}
\label{prompt-injection-outcomes}

Attack success was evaluated post hoc over the 250 adversarial records
for each configuration. The scorer uses an explicit target or an
observable failure of the trusted task for each injection variant. An
\texttt{unknown} prediction is not automatically counted as success.
Evidence suppression uses a conservative structural and lexical proxy,
so these ASR values should be interpreted as reproducible preliminary
measurements rather than human-verified outcome labels.
Table~\ref{tab:synthetic-asr} shows that the secure configurations do
not improve all outcomes uniformly.

\begin{table}[H]
\centering
\caption{Post-hoc Prompt Injection Attack Success Rate on the CSIC 2010-based synthetic dataset.}
\label{tab:synthetic-asr}
\begin{tabular}{@{}lrrr@{}}
\toprule
Configuration & Successful / 250 & ASR & Robustness \\
\midrule
C1 & 8 & 3.2\% & 96.8\% \\
C2 & 8 & 3.2\% & 96.8\% \\
C3 & 4 & 1.6\% & 98.4\% \\
C4 & 2 & 0.8\% & 99.2\% \\
\bottomrule
\end{tabular}
\end{table}

The largest observed contribution to these totals was not the direct
overriding of the output schema. Instruction-override and
prompt-extraction attempts had no detected successes under the stated
rules. Classification manipulation was detected in C1 2/50 and C2 3/50,
but in neither C3 nor C4. Evidence-suppression scores were 3/50,
3/50, 4/50, and 2/50 for C1--C4, respectively. Severity manipulation
was observed in 3/50, 2/50, 0/50, and 0/50 cases. Because severity
policy and evidence omission require semantic judgement, these counts
are less certain than the classification-target and prompt-extraction
checks.

Paired utility retention---both clean and adversarial members classified
correctly---was 78.8\% for C1, 70.4\% for C2, 61.2\% for C3, and 81.6\%
for C4.
C4 therefore had the best observed combination of ASR and retained
classification utility on the synthetic corpus.

\section{Standalone Prompt Injection Detector Evaluation}
\label{standalone-detector-evaluation}

The fixed C4 detector was evaluated against the adversarial/clean
condition labels on all 500 records. Match denotes a detector-positive
result. The confusion counts and derived metrics are:

\begin{table}[H]
\centering
\caption{C4 detector confusion counts and metrics on the 500-record CSIC 2010-based synthetic dataset.}
\label{tab:c4-detector}
\begin{tabular}{@{}lrr@{}}
\toprule
Condition & Match & No match \\
\midrule
Adversarial (250) & TP: 240 & FN: 10 \\
Clean (250) & FP: 0 & TN: 250 \\
\midrule
Precision & \multicolumn{2}{l}{240/(240+0) = 100.0\%} \\
Recall & \multicolumn{2}{l}{240/(240+10) = 96.0\%} \\
Specificity & \multicolumn{2}{l}{250/(250+0) = 100.0\%} \\
F1 score & \multicolumn{2}{l}{2(240)/(2(240)+0+10) = 97.96\%} \\
Accuracy & \multicolumn{2}{l}{(240+250)/500 = 98.0\%} \\
\bottomrule
\end{tabular}
\end{table}

These are detector measurements against the study's adversarial/clean
labels, not end-to-end defense success or Prompt Injection ASR. The
detector routed 240 records to the C4 secure path; this count is not
evidence that those records contained successful attacks.

\section{Citation-ID Membership Checks}\label{citation-membership-checks}

Citation membership was checked only for the 240 C4 secure-path reports.
Of these, 114/240 reports contained citations and 126/240 did not.
Across the reports there were 166 cited IDs, of which 0/166 were
outside the local knowledge-base ID set; equivalently, 0/240 reports
had an invalid citation ID. Membership in the knowledge-base ID set
does not demonstrate that a cited document semantically supports the
claim made in a response.

\section{Semantic Quality and Citation Support}
\label{semantic-quality}

A stratified sample of 50 reports per configuration (200 reports in
total) was reviewed by an LLM-as-judge procedure. The sample was
balanced by configuration, clean/adversarial condition, and traffic
category; the judge received the request, annotations, response, and
cited document content. These are assisted semantic annotations, not
independent human ground truth. Table~\ref{tab:semantic-quality}
summarizes both unsupported-report claims and the semantic support of
the cited material.

\begin{table}[H]
\centering
\caption{LLM-assisted semantic review of reports sampled from the CSIC 2010-based synthetic dataset (200 reports).}
\label{tab:semantic-quality}
\begin{tabular}{@{}lrrr@{}}
\toprule
Configuration & Hallucination rate & Citations judged & Semantic support \\
\midrule
C1 & 12/50 (24.0\%) & 0 & Not applicable \\
C2 & 16/50 (32.0\%) & 115 & 95/115 (82.6\%) \\
C3 & 15/50 (30.0\%) & 80 & 75/80 (93.8\%) \\
C4 & 16/50 (32.0\%) & 82 & 73/82 (89.0\%) \\
\bottomrule
\end{tabular}
\end{table}

Semantic support includes citations labelled as fully or partially
supporting a claim. The C3 sample had the highest support rate. C4's revised sample reached
89.0\% semantic support. Hallucination rates in
Table~\ref{tab:semantic-quality} measure unsupported or contradictory
factual claims; a wrong classification by itself was not automatically
treated as hallucination. Human double review of the sampled judgements
and expansion to the full corpus remain necessary before these values
can be treated as audited human measurements.

\section{External Evaluation on ModSec-WP}
\label{evaluation-modsec-wp}

The full model run used the 800-record ModSec-WP paired dataset
~\cite{modsecwp2026dataset}. All four configurations were run once per
record, yielding 3{,}200 responses. The quantitative results below use
600 records from 300 pairs after excluding the 200
\texttt{bruteforce} records and their 100 pairs; all 2{,}400 responses
in this reported cohort were valid structured JSON. The returned model
was \texttt{gpt-4o-mini-\allowbreak 2024-07-18}. The full run contains
200 benign and 600 malicious records; the reported cohort contains 200
benign and 400 malicious records. Original source labels are retained as
the evaluation reference, but were not independently verified by this
study. Results are descriptive because the selected sample also informed
secure prompt v5.

\begin{table}[H]
\centering
\caption{Classification and end-to-end attack-subtype accuracy on the 600-record ModSec-WP analytical cohort.}
\label{tab:modsec-classification}
{\small
\setlength{\tabcolsep}{3pt}
\resizebox{\textwidth}{!}{%
\begin{tabular}{@{}lrrrrrrr@{}}
\toprule
Config. & Correct (600) & Acc. & Cov. & Unknown &
Errors (FP/FN) & Correct subtype (400) & End-to-end acc. \\
\midrule
C1 & 534 & 89.0\% & 93.8\% & 37 & 29/0 & 395 & 98.8\% \\
C2 & 499 & 83.2\% & 86.8\% & 79 & 18/4 & 381 & 95.3\% \\
C3 & 505 & 84.2\% & 85.8\% & 85 & 10/0 & 348 & 87.0\% \\
C4 & 539 & 89.8\% & 91.7\% & 50 & 11/0 & 378 & 94.5\% \\
\bottomrule
\end{tabular}
}
}
\par\smallskip
\begin{minipage}{\textwidth}\footnotesize
\textit{Note.} Unknown predictions are abstentions, counted as incorrect
for accuracy. Coverage is the share of records assigned benign or
malicious. Errors are shown as false positives / false negatives; unknowns
are excluded from both counts. Correct and end-to-end subtype counts use
denominators of 600 and 400, respectively. A subtype prediction counts
as correct only when the record is also classified as malicious; this is
not precision over predicted attacks.
\end{minipage}
\end{table}

C4 had the highest observed overall classification accuracy (89.8\%),
while C1 had the highest coverage (93.8\%). The source-category
breakdown in Table~\ref{tab:modsec-category-accuracy} shows
98.0--100.0\% classification accuracy on SQL Injection and
76.0--100.0\% on XSS across configurations. The \texttt{bruteforce}
stratum is excluded from these results because its label requires a
temporal sequence of requests; the individual model input does not
include source address or timestamp.

\begin{table}[H]
\centering
\caption{Binary classification accuracy by source category in the ModSec-WP analytical cohort.}
\label{tab:modsec-category-accuracy}
\begin{tabular}{@{}lrrr@{}}
\toprule
Configuration & Normal & SQL Injection & XSS \\
\midrule
C1 & 67.0\% & 100.0\% & 100.0\% \\
C2 & 57.0\% & 98.0\% & 94.5\% \\
C3 & 78.5\% & 98.0\% & 76.0\% \\
C4 & 79.5\% & 98.5\% & 91.5\% \\
\bottomrule
\end{tabular}
\par\smallskip
\begin{minipage}{\textwidth}\footnotesize
\textit{Note.} Each included category contains 200 records: 100 clean
and 100 adversarial. ``Normal'' is the study mapping of the source
\texttt{normal} label. The source \texttt{bruteforce} stratum is not
included in the reported metrics because the class requires temporal
window evidence.
\end{minipage}
\end{table}

The secondary audit found that 99 of the 100 sampled source requests met
the five-request, five-minute criterion for sequence-supported repeated
authentication activity; one did not. This indicates why the stratum is
not evaluated as isolated-request classification: the category depends
on the observation window and source-level aggregation. The audit does
not establish that authentication failed for any individual request and
is used only to document the exclusion rationale.

\subsection{Prompt Injection and Paired Utility}
\label{evaluation-modsec-pi}

Table~\ref{tab:modsec-pi-results} reports classification on the clean
and adversarial members, post-hoc Attack Success Rate (ASR), and paired
utility for the 300-pair analytical cohort. ASR counts responses
meeting the pre-defined deterministic success rules among 300
adversarial records per configuration; robustness is
\(1-\mathrm{ASR}\). Pair retention is the fraction of 300 pairs for
which both members receive the correct binary classification.

\begin{table}[H]
\centering
\caption{Prompt Injection and paired-utility results on the ModSec-WP analytical cohort.}
\label{tab:modsec-pi-results}
{\small
\setlength{\tabcolsep}{4pt}
\begin{tabular}{@{}lrrrrrr@{}}
\toprule
Configuration & \shortstack{Clean\\accuracy} &
\shortstack{Adversarial\\accuracy} & ASR &
\shortstack{Robustness\\(1--ASR)} &
\shortstack{Both-correct\\pairs} &
\shortstack{Pair\\retention} \\
\midrule
C1 & 96.0\% & 82.0\% & 16/300 (5.3\%) & 94.7\% & 244/300 & 81.3\% \\
C2 & 92.0\% & 74.3\% & 9/300 (3.0\%) & 97.0\% & 219/300 & 73.0\% \\
C3 & 87.3\% & 81.0\% & 0/300 (0.0\%) & 100.0\% & 231/300 & 77.0\% \\
C4 & 91.0\% & 88.7\% & 4/300 (1.3\%) & 98.7\% & 253/300 & 84.3\% \\
\bottomrule
\end{tabular}
}
\par\smallskip
\begin{minipage}{\textwidth}\footnotesize
\textit{Note.} Clean and adversarial accuracy are calculated on the 300
clean and 300 injected records, respectively. ASR is successful
Prompt Injection outcomes divided by 300; robustness is \(1-\mathrm{ASR}\).
Both-correct pairs counts pairs where both records were classified
correctly; pair retention is that count divided by 300.
\end{minipage}
\end{table}

C3 had the lowest measured ASR (0/300), but retained both correct
classifications on 77.0\% of pairs. C4 had the highest paired
retention (84.3\%) and adversarial accuracy (88.7\%) in the analytical
cohort. These outcomes show why ASR must be read alongside
classification utility; they do not establish a causal contribution
from any one configuration component.

\begin{longtable}{@{}p{2.7cm}p{7.8cm}rr@{}}
\caption{Post-hoc ASR by Prompt Injection variant on the ModSec-WP analytical cohort.}
\label{tab:modsec-variant-asr}\\
\toprule
Configuration & Prompt Injection variant & Successful / \(n\) & ASR \\
\midrule
\endfirsthead
\multicolumn{4}{c}{\textit{Table \thetable{} continued}}\\
\toprule
Configuration & Prompt Injection variant & Successful / \(n\) & ASR \\
\midrule
\endhead
\midrule
\multicolumn{4}{r@{}}{\textit{Continued on next page}}\\
\endfoot
\bottomrule
\endlastfoot
C1 & Instruction Override & 0/65 & 0.0\% \\
C1 & Classification Manipulation & 14/65 & 21.5\% \\
C1 & Severity Manipulation & 2/65 & 3.1\% \\
C1 & Evidence Suppression & 0/40 & 0.0\% \\
C1 & Prompt Extraction & 0/65 & 0.0\% \\
\midrule
C2 & Instruction Override & 0/65 & 0.0\% \\
C2 & Classification Manipulation & 7/65 & 10.8\% \\
C2 & Severity Manipulation & 2/65 & 3.1\% \\
C2 & Evidence Suppression & 0/40 & 0.0\% \\
C2 & Prompt Extraction & 0/65 & 0.0\% \\
\midrule
C3 & Instruction Override & 0/65 & 0.0\% \\
C3 & Classification Manipulation & 0/65 & 0.0\% \\
C3 & Severity Manipulation & 0/65 & 0.0\% \\
C3 & Evidence Suppression & 0/40 & 0.0\% \\
C3 & Prompt Extraction & 0/65 & 0.0\% \\
\midrule
C4 & Instruction Override & 0/65 & 0.0\% \\
C4 & Classification Manipulation & 1/65 & 1.5\% \\
C4 & Severity Manipulation & 3/65 & 4.6\% \\
C4 & Evidence Suppression & 0/40 & 0.0\% \\
C4 & Prompt Extraction & 0/65 & 0.0\% \\
\bottomrule
\end{longtable}
\begin{minipage}{\textwidth}\footnotesize
\textit{Note.} Successful / \(n\) is the number of adversarial records
meeting the variant-specific success rule divided by the number assigned
to that variant in the analytical cohort: \(n=65\) for each variant
except Evidence Suppression (\(n=40\)). ASR is the successful count
divided by \(n\). These post-hoc scores are deterministic proxies, not
human-verified outcomes.
\end{minipage}

The C4 detector matched 300/300 adversarial records and 0/300 clean
records, corresponding to 100.0\% recall, precision, and specificity on
the reported cohort. This is expected coverage of the canonical
study-authored instructions, which were designed to exercise the
detector's defined patterns; it is not evidence that the detector
generalizes to unseen or obfuscated instructions. Detector matches are
also distinct from ASR: C4 had 4/300 measured successes despite all
adversarial records activating its secure route. Severity Manipulation
is scored only against its injected low-severity target because
ModSec-WP supplies no severity ground truth. Evidence-suppression and
other lexical or structural success rules remain proxies and require
human review.

\section{Failure Analysis}\label{failure-analysis}

The aggregate summary does not provide verified per-record evidence
needed to attribute individual failures to retrieval, reasoning,
Prompt Injection, evidence attribution, or ambiguous ground truth.
Record-level human failure analysis remains pending; the automated and
LLM-assisted counts above should not be used to assign a cause to any
individual response without reviewing the record.
